\section{Estrutura Analítica da Obra}

Para prover uma fundação científica e epistêmica condizente com a rigidez acadêmica de nível \textit{Qualis A1}, este manuscrito está meticulosamente estruturado em quatro eixos centrais, garantindo fluidez lógica da física fundamental à aplicação clínica e regulação estatal:

\begin{enumerate}
    \item \textbf{Parte I --- Fundamentos Teóricos e Biomecânicos (Caps. 1 a 3)}: Descreve a arquitetura matemática do DT, as ontologias computacionais e os princípios basilares da modelagem preditiva aplicados à medicina personalizada de precisão. O rigor da transição dos dados fisiológicos brutos para as matrizes de elementos finitos é abordado detalhadamente.
    \item \textbf{Parte II --- Aplicações dos Gêmeos Digitais na Odontologia (Caps. 4 a 8)}: Consolida a práxis clínica suportada por evidências. Aprofunda-se na análise vetorial e cinemática da Ortodontia (em especial a terapia com alinhadores invisíveis)~\cite{li2025impacts, olawade2026advancing}, o planejamento estocástico na Implantodontia para maximização de contato osso-implante~\cite{alshadidi2024surface}, e as fronteiras da Endodontia preditiva no comportamento de ligas metálicas no interior dos canais radiculares~\cite{elsaid2025digital}.
    \item \textbf{Parte III --- Ecossistemas de IA e Orquestração Aberta (Caps. 9 a 12)}: Apresenta o detalhamento técnico do \textit{OpenCode Ecosystem}. Explora os pipelines de \textit{Machine Learning} em aquisição volumétrica (CBCT) e demonstra como enxames de agentes de IA operam de forma autônoma para a segmentação de imagens e tomada de decisão preditiva~\cite{opencode2026ecosystem, mdpi2025digital}.
    \item \textbf{Parte IV --- Impacto Social, Equidade e Marco Regulatório (Caps. 13 a 16)}: Um ensaio crítico sobre as implicações bioéticas do determinismo preditivo. Aborda a cibersegurança e o sigilo de dados genéticos~\cite{springer2026realising}, concluindo com o papel transformador da odontologia in-silico como ferramenta de mitigação da desigualdade nas políticas públicas e no seio do Sistema Único de Saúde (SUS)~\cite{cuadros2023assessing}.
\end{enumerate}

\destaque{Esta obra não visa apenas descrever a tecnologia existente, mas estabelecer o protocolo-base de como os Gêmeos Digitais devem ser construídos, validados e integrados no cotidiano da Odontologia baseada em dados.}
